\begin{description}
\item[(T11.1a)]\mbox{}
  \begin{itemize}
  \item $X$ and $l$ increase with $r$, while $T$, $\epsilon_{\mathrm{nuc}}$,
    and $Y$ decrease with $r$. So A and B must be $X$ and $l$.
  \item For a $1\,\mathrm{M_\odot}$ star of age 4\,Gyr, $X \neq 0$ at centre.
    So B cannot be $X$.
  \item Also, $l = 0$ at centre. Matches with B.
  \item Both $X$ and $l$ must saturate to their maximum values just outside
    the H-burning core, and this happens at $r/R = 0.3$ here.
  \item Each of $T$, $\epsilon_{\mathrm{nuc}}$, and $Y$ have their maximum
    value at the centre and decrease outwards. So C is one of these.
  \item But $Y$ cannot be (close to) zero at the surface ($r/R = 1$) compared
    to its central value.
  \item $\epsilon_{\mathrm{nuc}}$ becomes extremely small just outside the
    core ($r/R \lesssim 0.3$), and not at the surface.
  \item Therefore, C must be $T(r)$ (the surface value is
    ${\sim}\,10^3$\,K, compared to the maximum value at the centre
    ${\sim}\,10^6$\,K).
  \end{itemize}
  Therefore,
  \begin{align*}
    \mathrm{A} &\longrightarrow X(r) \tag*{(2.0)}\\
    \mathrm{B} &\longrightarrow l(r) \tag*{(2.0)}\\
    \mathrm{C} &\longrightarrow T(r) \tag*{(2.0)}
  \end{align*}

\item[(T11.1b)] Maximum value of $X$ occurs outside the core, and is given by
  \[
    X_{\mathrm{s}} = 1 - Y_{\mathrm{s}} - Z_{\mathrm{s}}
    = 1 - 0.28 - 0.02 = 0.70
  \]
  \hfill(1.0)

  Note $Z$ remains constant throughout the star in the absence of
  gravitational settling and diffusion (given).

  From the given graph, the central hydrogen mass fraction, $X_{\mathrm{c}}$,
  is approximately
  \[
    X_{\mathrm{c}} \approx 0.55 X_{\mathrm{s}} = 0.55 \times 0.70 = 0.39
  \]
  \hfill(1.5)

  Then, the helium mass fraction at the centre is
  \[
    Y_{\mathrm{c}} = 1 - X_{\mathrm{c}} - Z_{\mathrm{c}}
    = 1 - 0.39 - 0.02 = 0.59
  \]
  \hfill(0.5)

\item[(T11.1c)]\mbox{}
  \begin{itemize}
  \item The curve for $\epsilon_{\mathrm{nuc}}$ must mirror that of $l$: any
    reasonable smooth curve between $(0,1)$ and $(0.3,0)$. \hfill(2.5)
  \item The curve for $Y$ must mirror that of $X$, and flatten outside the
    core to the normalized value corresponding to $Y_{\mathrm{s}} = 0.28$.
    Previous result of $Y_{\mathrm{c}} = 0.59$ gives this normalized value:
    $0.28/0.59 \approx 0.47$. \hfill(2.5)
  \end{itemize}
  % FIGURE NEEDED: answer graph with hand-drawn Y and \epsilon_{nuc} curves
  % overlaid on the A/B/C plot (2025 theory solutions, page 44 of 59)

\item[(T11.2a)] During the main sequence phase, the mass fraction of helium at
  the centre increases steadily, up to a maximum value of
  $Y_{\mathrm{c,\,max}} = 1 - Z_{\mathrm{c}}$. Here that maximum value is
  reached at an age of around 8.5\,Gyr, from the plot of $Y_{\mathrm{c}}$ vs
  age.
  \[
    t_{\mathrm{MS}} = 8.5 \times 10^{9}\,\mathrm{yr}
  \]
  \hfill(1.0)

\item[(T11.2b)] During the helium burning phase, the mass fraction of helium
  at the centre decreases steadily, down to 0. The inset in the plot of
  $Y_{\mathrm{c}}$ vs age shows that this corresponds to ages between
  $11.896 \times 10^{9}$\,yr and $11.960 \times 10^{9}$\,yr, corresponding to
  a duration of 64\,Myr.
  \[
    \Delta t_{\mathrm{He}} = 64 \times 10^{6}\,\mathrm{yr}
  \]
  \hfill(1.0)

\item[(T11.2c)] Read-off from the plot of $\log L/L_\odot$ vs age, we get age
  ${\approx}\,4 \times 10^{9}$\,yr. \hfill(1.0)

  At this age, read-off from plot of $Y_{\mathrm{c}}$ vs age, we get
  \[
    Y_{\mathrm{c}} \approx 0.60 \implies
    X_{\mathrm{c}} \approx 1 - 0.60 - 0.02 = 0.38
  \]
  \hfill(1.0)

  Therefore the fraction of original H burnt is
  \[
    f_{\mathrm{H}} = 1 - \frac{X_{\mathrm{c}}}{X_0}
    = 1 - \frac{X_{\mathrm{c}}}{1 - Y_0 - Z_0} = 0.46
  \]
  \hfill(1.0)

\item[(T11.2d)]
  \[
    X_0 = 1 - Y_0 - Z_0 = 1 - 0.28 - 0.02 = 0.70
  \]
  At this stage 60\% of H has been burnt at the centre. So
  $X_{\mathrm{c}} = 0.40 X_0 = 0.28$. Therefore,
  \[
    Y_{\mathrm{c}} = 1 - X_{\mathrm{c}} - Z_{\mathrm{c}}
    = 1 - 0.28 - 0.02 = 0.70
  \]
  \hfill(2.0)

  From the graph of $Y_{\mathrm{c}}$ vs age, this corresponds to an age of
  approximately $5 \times 10^{9}$\,yr. At this age, the value of radius is
  approximately $1\,\mathrm{R_\odot}$ ($\log R/R_\odot \approx 0$), from the
  graph of $\log R/R_\odot$ vs age.
  \[
    R_1 = 1\,\mathrm{R_\odot}
  \]
  \hfill(1.0)

\item[(T11.2e)] For P:
  $\log T_{\mathrm{eff}} = 3.74 \implies T_{\mathrm{eff}} \approx 5500$\,K
  (the $L$, being flat, does not help here). \hfill(1.0)

  Read-off from plot of $T_{\mathrm{eff}}$ vs age, we get age
  ${\approx}\,11 \times 10^{9}$\,yr. At this age, read-off from plot of
  $\log R/R_\odot$ vs age, we get $\log R/R_\odot \approx 0.25 \implies
  R_{\mathrm{P}} = 1.78\,\mathrm{R_\odot}$. \hfill(1.0)

  For Q: $\log L/L_\odot = 2.5$. (One can also use
  $\log T_{\mathrm{eff}} = 3.60 \implies T_{\mathrm{eff}} \approx 3981$\,K
  which would require rounding off to 3950\,K at the next step.) \hfill(1.0)

  Read-off from (the inset) plot of $\log L/L_\odot$ vs age, we get age
  ${\approx}\,11.884 \times 10^{9}$\,yr. Notice that this value of
  $\log L/L_\odot$ occurs twice, and one has to take the first one, i.e., the
  ascending RGB branch.

  At this age, read-off from (inset) plot of $\log R/R_\odot$ vs age, we get
  $\log R/R_\odot \approx 1.6 \implies R_{\mathrm{Q}} = 39.81\,\mathrm{R_\odot}$.
  \hfill(1.0)

\item[(T11.3a)]
  \begin{align*}
    \frac{dq}{dx} &= \frac{dq}{dm}\frac{dm}{dr}\frac{dr}{dx}
      \tag*{(0.5)}\\
      &= \frac{1}{M}\left[4\pi (Rx)^2 \rho\right] R\\
      &= 3\left(\frac{4\pi R^3}{3M}\right) x^2 \rho
       = 3 \frac{1}{\bar{\rho}} x^2 \rho
      \tag*{(1.0)}\\
    \therefore\ \frac{dq}{dx} &= 3 x^2 \sigma
      \tag*{(0.5)}
  \end{align*}

\item[(T11.3b)] Let us denote the line as
  \[
    \log\sigma = \alpha \log x + \beta
  \]
  From the graph, we use the two end points $(-0.5, 1.0)$ and $(-0.1, -1.0)$.
  Therefore,
  \begin{align*}
    1.0 &= -0.5\alpha + \beta\\
    -1.0 &= -0.1\alpha + \beta
  \end{align*}
  Solving, we get $\alpha = -5$ and $\beta = -1.5$. Therefore,
  \begin{align*}
    \log\sigma &= -5\log x - 1.5
      &&\text{for } -0.5 < \log x < -0.1
      \tag*{(2.0)}\\
    \therefore\ \sigma(x) &= 10^{-1.5} x^{-5}
      &&\text{for } 0.32 < x < 0.80\\
    \implies \sigma(x) &= 0.030\,x^{-5}
      &&\text{for } 0.32 < x < 0.80
      \tag*{(2.0)}
  \end{align*}

\item[(T11.3c)]
  \[
    \frac{dq}{dx} = 3x^2\sigma = 3x^2(0.030\,x^{-5}) = 0.090\,x^{-3}
  \]
  \hfill(1.0)

  Integrating, we get
  \begin{align*}
    q(x) &= \int 0.090\,x^{-3}\,dx
      \tag*{(1.0)}\\
      &= -\frac{0.045}{x^2} + C
      \qquad (C \text{ is a constant of integration})
      \tag*{(1.0)}
  \end{align*}
  Now, given that $q(x = 0.35) = 0.7$. Using this, we get $C = 1.067$.
  \hfill(2.0)

  Therefore (rounding off to within 0.005),
  \[
    q(x) = 1.065 - \frac{0.045}{x^2}
    \qquad \text{for } 0.32 < x < 0.80
  \]
  \hfill(1.0)

\item[(T11.3d)]
  \[
    \frac{dq}{dx} = 3x^2\sigma = 3x^2(Ax + B) = 3Ax^3 + 3Bx^2
  \]
  \hfill(1.0)

  Integrating, we get
  \begin{align*}
    q(x) &= \int (3Ax^3 + 3Bx^2)\,dx
      \tag*{(1.0)}\\
      &= \frac{3A}{4}x^4 + Bx^3 + D
      \qquad (D \text{ is a constant of integration})
      \tag*{(1.0)}
  \end{align*}
  Three boundary conditions applied sequentially to find the constants:
  \begin{itemize}
  \item Mass at the centre is zero, i.e., $q(0) = 0 \implies D = 0$
    \hfill(0.5)
  \item $\rho(0) = 80\bar{\rho}$ (given), i.e.,
    $\sigma(0) = 80 \implies B = 80$ \hfill(1.5)
  \item Given $q(0.25) = 0.50$. Therefore,
    \begin{align*}
      0.5 &= \frac{3A}{4}(0.25)^4 + 80\,(0.25)^3\\
      \implies A &= \frac{4}{3}\,\frac{0.50 - 80(0.25)^3}{(0.25)^4}\\
      \implies A &= -256
      \tag*{(2.0)}
    \end{align*}
  \end{itemize}
  Therefore, putting the values of $A$ and $B$,
  \[
    q(x) = -192 x^4 + 80 x^3 \qquad \text{for } 0 \le x \le 0.32
  \]
  \hfill(1.0)

\item[(T11.3e)] The sketch should be a smooth curve for $q(x)$ from $(0,0)$
  to $(1,1)$, with zero slope at both ends, passing through the two given
  points $(0.25, 0.5)$ and $(0.35, 0.7)$; the initial part looks like a
  cubic--quartic, and the latter part looks like a negative inverse-square
  with a dc shift.
  % FIGURE NEEDED: hand-drawn answer sketch of q(x) vs x through (0.25,0.5)
  % and (0.35,0.7) (2025 theory solutions, page 50 of 59)
\end{description}
